\section{Direct Rebound Effects in AI Systems}
\label{sec:direct}

The most immediate form of rebound effect in the context of AI 
concerns the infrastructure required to develop and operate AI 
systems. While individual AI models have become increasingly 
efficient in terms of energy consumption per query, this progress 
has not translated into a reduction in absolute energy use. 
Instead, lower costs per operation have stimulated a dramatic 
increase in demand, a pattern consistent with the Jevons Paradox 
described in the preceding section.

A concrete illustration of this dynamic is the development of 
graphics processing units (GPUs), the primary hardware used for 
training and deploying AI models. Despite continuous improvements 
in hardware efficiency, NVIDIA shipped 3.7 million data center 
GPUs in 2024 -- more than one million more units than in the 
previous year -- driven by surging demand from AI applications 
\cite{luccioni2025}. The data centers housing this hardware have 
similarly improved their Power Usage Effectiveness (PUE), yet 
their absolute energy consumption has continued to rise 
\cite{luccioni2025}. This illustrates the core mechanism of the 
direct rebound effect: efficiency gains at the component level 
are offset by growth in the overall scale of deployment.

The scale of this growth is substantial. According to projections 
by the International Energy Agency, global data center electricity 
consumption is expected to reach between 650 and 1,050 TWh by 
2026, driven primarily by the rapid expansion of AI infrastructure 
\cite{iea2025ai}. To put this figure in context, the lower bound 
alone would exceed the total annual electricity consumption of 
several mid-sized European countries.

The case of Google DeepMind illustrates the paradox at the 
facility level. The deployment of an AI system to optimize 
cooling in Google's data centers reduced cooling energy 
consumption by up to 40 percent \cite{mhlanga2025}. While this 
represents a significant efficiency gain, the broader trend 
points in the opposite direction: the same technology companies 
achieving these efficiency improvements are simultaneously 
expanding their data center capacity at an unprecedented rate. 
Efficiency, in this context, appears to function as an enabler 
of growth rather than a substitute for it.

A further challenge in assessing the direct environmental impact 
of AI is the lack of standardized reporting. Technology companies 
commonly disclose only a narrow range of environmental metrics, 
which impedes comprehensive lifecycle assessments and leaves the 
full extent of AI's resource consumption largely opaque 
\cite{luccioni2025}. This absence of transparent data not only 
complicates academic analysis but also undermines the capacity 
of policymakers to design effective regulatory responses.